\documentclass[10pt,a4paper]{amsart}
\usepackage[margin=19mm]{geometry}
\usepackage{amsmath,amssymb,mathtools}
\usepackage[scr=rsfs,cal=euler]{mathalfa}
\usepackage{xfrac}
\usepackage[unicode,hidelinks,bookmarks=false]{hyperref}
\newcommand{\ie}{i.e.\ }
\newcommand{\cf}{cf.\ }
\newcommand{\resp}{respectively\ }
\newcommand{\Subsection}{\subsection}
\providecommand{\nobreakdash}{\nobreak\hbox{-}}
\setlength{\emergencystretch}{2em}
\multlinegap0pt
\usepackage{tikz}
\usepackage{mathtools}
\usepackage{bm}
\newtheorem{theorem}{Theorem}
\title{A central limit theorem\\ for the signatures of 2-bridge knots}
\author{Cody Baker, Moshe Cohen, Henry Dam, Rebecca Felber, Neal Madras, Ritvik Saha, Daisy Thackrah}
\date{Source: arXiv:2604.21107v1 (CC BY 4.0)}
\begin{document}
\maketitle

\noindent\emph{Source and licence.} A central limit theorem\\ for the signatures of 2-bridge knots. Version arXiv:2604.21107v1 (CC BY 4.0), distributed by arXiv under the Creative Commons Attribution 4.0 International licence, http://creativecommons.org/licenses/by/4.0/, as recorded in the arXiv licence metadata for this exact version and shown on the abstract page https://arxiv.org/abs/2604.21107v1. Full licence notice: \texttt{licenses/arxiv-2604.21107-CC-BY-4.0.txt}.\par\smallskip

\noindent\textit{Editorial scope.} Counted unit: the article's theorem thm:CLT (source0.tex lines 262--273). The article's complete original proof of it is delivered with it (source0.tex lines 854--873, one \texttt{proof} environment of 20 lines, which the article places away from the statement). No mathematics is written, repaired or re-derived: the statement and the proof are the article's own lines, in the article's own notation, and no retained line is substituted or re-flowed.  Retained uncounted as the source-local context the proof needs: lines 653--656; lines 838--841; lines 845--848.  Nothing was omitted from any retained span, so the package carries no omission record.  No retained line contains a citation, so the wrapper carries no bibliography and no bibliography entry.  No retained line contains a figure environment, an image-inclusion command or drawing code, so no image is delivered and the package declares no asset dependency.  Editorial formatting only, in addition: an AMS \texttt{amsart} preamble that restates the article's own declarations and macros used by the retained lines, namely the environments \texttt{theorem} on separate counters, together with the standard packages those lines need.  The retained environments print in this document as Theorem 1 in order of appearance, whereas the article prints the same items with the numbering of its own full document.  The complete native source of the article as distributed in the arXiv e-print for this exact version is bundled unchanged as \texttt{sources/198999/source0.tex}.
\begin{equation}
    \label{eq.CLT2m}
    \lim_{c\rightarrow\infty}  \frac{s\left(c, \,\leq t\sqrt{c}\right)}{|T(c)|}   \;=\; \Phi(t).
\end{equation}

\begin{equation*}
   \label{eq.KTT}
   2\, |K(c)|  \;=\;  |T(c)| \,+\,  |T_p(c)|    \hspace{5mm}\hbox{and}\hspace{5mm}   k(c,\sigma)  \;=\; \frac{1}{2}\left(s(c,\sigma)+s_p(c,\sigma)\right).
\end{equation*}

\begin{equation}
   \label{eq.KTlim}
   \lim_{c\rightarrow\infty}  \frac{|T(c)|}{2|K(c)|}  \;=\; 1.
\end{equation}

\begin{theorem}
\label{thm:CLT}
For $c\geq 3$, let $K(c)$ be the set of 2-bridge knots of a given crossing number $c$ with each knot counted exactly once.  Let $k(c,\leq x)$ be the number of knots in $K(c)$ with signature at most $x$.  
For every real number $t$,
\begin{equation}
    \label{eq.CLTk}
    \lim_{c\rightarrow\infty}  \frac{k\left(c, \,\leq t\sqrt{c}\right)}{|K(c)|}   \;=\; \Phi(t),
\end{equation}
where $\Phi$ is the cumulative distribution  function of the standard normal distribution. 
That is, over the set $K(c)$ of 2-bridge knots with crossing number $c$, the distribution of signatures 
is asymptotically normally distributed  with mean 0 and variance $c$ as $c$ tends to infinity.
\end{theorem}

\begin{proof}[Proof of Theorem \ref{thm:CLT}]
Following Equation \eqref{eq.KTT} we have
\begin{eqnarray*}
     \frac{k(c,\leq\sigma)}{|K(c)|}  & = & \frac{s(c,\leq \sigma)}{2\,|K(c)|}  \,+\, \frac{s_p(c,\leq \sigma)}{2\,|K(c)|} 
     \\
     & = &    \frac{s(c,\leq \sigma)}{|T(c)|} \,\frac{|T(c)|}{2\,|K(c)|}  \,+\, \frac{s_p(c,\leq \sigma)}{2\,|K(c)|}  \,.
\end{eqnarray*}
It follows that 
\begin{eqnarray}
   \nonumber
   \left|   \frac{k(c,\leq\sigma)}{|K(c)|}  \,-\,   \frac{s(c,\leq \sigma)}{|T(c)|} \right|   & \leq &        \frac{s(c,\leq \sigma)}{|T(c)|} \,
    \left|  \frac{|T(c)|}{2\,|K(c)|}   \,-\,1 \right|   \,+\,   \frac{|T_p(c)|}{2\,|K(c)|}
    \\
    \label{eq.knears}
       & \leq &       \left|  \frac{|T(c)|}{2\,|K(c)|}   \,-\,1 \right|   \,+\,   \left| \frac{2\,|K(c)|-|T(c)|}{2\,|K(c)|} \right|\,.
 \end{eqnarray}
By Equation \eqref{eq.KTlim}, the right-hand side of Equation \eqref{eq.knears} converges to 0 as $c\rightarrow\infty$, uniformly in $\sigma$.   
Therefore Equation \eqref{eq.CLTk} follows from 
Equation \eqref{eq.CLT2m} on the Central Limit Theorem for the set $T(c)$, and the theorem is proved.
\end{proof}

\end{document}
